
\documentclass{article}
\usepackage{colortbl}
\usepackage{makecell}
\usepackage{multirow}
\usepackage{supertabular}

\begin{document}

\newcounter{utterance}

\centering \large Interaction Transcript for game `bbh', experiment `snarks', episode 139 with reflexion{-}Qwen3.5{-}27B.
\vspace{24pt}

{ \footnotesize  \setcounter{utterance}{1}
\setlength{\tabcolsep}{0pt}
\begin{supertabular}{c@{$\;$}|p{.15\linewidth}@{}p{.15\linewidth}p{.15\linewidth}p{.15\linewidth}p{.15\linewidth}p{.15\linewidth}}
    \# & \multicolumn{2}{c}{Player} && \multicolumn{2}{c}{Game Master} \\
    \hline

    \theutterance \stepcounter{utterance}  
    & & & \multicolumn{4}{p{0.6\linewidth}}{
        \cellcolor[rgb]{0.9,0.9,0.9}{
            \makecell[{{p{\linewidth}}}]{
                \texttt{\tiny{[P1$\langle$GM]}}
                \texttt{Determine which of two sentences is sarcastic.} \\
\\ 
\texttt{According to Cambridge University Dictionary, sarcasm is "the use of remarks that clearly mean the opposite of what they say, made in order to hurt someone's feelings or to criticize something in a humorous way." Sarcastic sentences often contain satirical or ironic utterances, hyperboles, ambivalent or witty remarks.} \\
\\ 
\texttt{Q: Which statement is sarcastic?} \\
\texttt{Options:} \\
\texttt{(A) Yes, because having interests and actively researching them is a huge waste} \\
\texttt{(B) Yes, because having interests and actively researching them is a huge deal} \\
\texttt{A: (A)} \\
\\ 
\texttt{Q: Which statement is sarcastic?} \\
\texttt{Options:} \\
\texttt{(A) No one is going to disagree with you on this. Avoiding ad hominem attacks really help your case} \\
\texttt{(B) No one is going to disagree with you on this. Ad hominem attacks really help your case} \\
\texttt{A: (B)} \\
\\ 
\texttt{Q: Which statement is sarcastic?} \\
\texttt{Options:} \\
\texttt{(A) Consistency in the league's punishments? What do you think this is supposed to be, politics?} \\
\texttt{(B) Consistency in the league's punishments? What do you think this is supposed to be, moral?} \\
\texttt{A: (A)} \\
\\ 
\texttt{Q: Which statement is sarcastic?} \\
\texttt{Options:} \\
\texttt{(A) Thank you stranger on the internet, I won't trust your three word post without any sources over someone offering solid advice claiming to actually be from the country in question} \\
\texttt{(B) Thank you stranger on the internet, I will trust your three word post without any sources over someone offering solid advice claiming to actually be from the country in question} \\
\texttt{A:} \\
            }
        }
    }
     \\ \\

    \theutterance \stepcounter{utterance}  
    & \multicolumn{4}{p{0.6\linewidth}}{
        \cellcolor[rgb]{0.561,0.737,0.561}{
            \makecell[{{p{\linewidth}}}]{
                \texttt{\tiny{[P1$\rangle$GM]}}
                \texttt{A: (B)} \\
\\ 
\texttt{**Reasoning:**} \\
\texttt{Sarcasm involves saying the opposite of what is meant, often to criticize or mock.} \\
\texttt{*   **Option (A)** states a logical, sincere position: refusing to trust an unverified, short post over solid advice from an expert. This is a genuine statement of reasoning.} \\
\texttt{*   **Option (B)** claims the speaker *will* trust the unverified, three{-}word post over the expert advice. Since this is an irrational and illogical action, the statement is clearly meant to be ironic. The speaker is mocking the idea of trusting such low{-}quality information, implying that doing so would be foolish. Therefore, (B) is the sarcastic statement.} \\
            }
        }
    }
    & & \\ \\

    \theutterance \stepcounter{utterance}  
    & & & \multicolumn{2}{p{0.3\linewidth}}{
        \cellcolor[rgb]{0.9,0.9,0.9}{
            \makecell[{{p{\linewidth}}}]{
                \texttt{\tiny{[GM$|$GM]}}
                \texttt{Target: (B)} \\
            }
        }
    }
    & & \\ \\

    \theutterance \stepcounter{utterance}  
    & & & \multicolumn{2}{p{0.3\linewidth}}{
        \cellcolor[rgb]{0.9,0.9,0.9}{
            \makecell[{{p{\linewidth}}}]{
                \texttt{\tiny{[GM$|$GM]}}
                \texttt{game\_result = LOSE} \\
            }
        }
    }
    & & \\ \\

    \theutterance \stepcounter{utterance}  
    & & \multicolumn{4}{p{0.6\linewidth}}{
        \cellcolor[rgb]{0.95,0.95,0.95}{
            \makecell[{{p{\linewidth}}}]{
                \texttt{\tiny{[P1$\leftrightarrow$P]}}
                \texttt{Post{-}episode Reflexion (not sent to GM):} \\
\texttt{New plan: I incorrectly identified the sarcastic statement by assuming that the option containing an illogical or irrational action (trusting a bad source) was the sarcastic one. However, in the provided examples, sarcasm was defined by the use of irony to criticize a situation, often by equating a negative reality with something absurd (like politics) or by stating a positive{-}sounding phrase that clearly means the opposite in context. In the final question, Option (A) uses the phrase "Thank you stranger on the internet" to mock the person, followed by a sincere rejection of their advice; the sarcasm lies in the polite address used to deliver a harsh critique. Option (B) is a literal statement of a foolish action, which lacks the ironic "opposite meaning" structure required by the definition. Next time, I will look for the option that uses polite or positive language to deliver a negative critique, or equates a serious issue with an absurd concept, rather than simply selecting the option that describes an illogical behavior. I will specifically analyze the tone of the opening phrase and the relationship between the premise and the conclusion to identify the ironic contrast.} \\
            }
        }
    }
    & \\ \\

\end{supertabular}
}

\end{document}
